% Appendix B - Detailed requirements analysis.
% Source: docs/requirements.md (transcription into LaTeX).
% When docs/requirements.md changes, propagate it here. Traceability
% is currently handled manually.

\umaappendix{Detailed requirements analysis}
\label{ap:requirements}

This appendix reproduces the requirements analysis of UnivOrch v1
that was developed in Phase~2 of the project. The structured
summary of the same analysis is in \Cref{sec:analisis} of
\Cref{ch:desarrollo}. The authoritative source lives at
\texttt{docs/requirements.md} in the project repository; if the two
diverge in the future, that document takes precedence.

\section{Purpose and scope}

The analysis covers the generic orchestrator core and the teaching
application built on top. It records use cases, functional
requirements, non-functional requirements, constraints and
explicitly out-of-scope items. Architecture and technology choices
do not appear here; they are deferred to \Cref{ch:desarrollo} and
\Cref{ch:estado-arte} respectively.

\section{Actors}

\Cref{tab:actors-detail} lists the actors of the system. A single
person can hold different roles in different branches of the tree;
the role lives on the folder, not on the user record (DEC-021).

\begin{table}[ht]
  \centering
  \small
  \begin{tabular}{p{0.22\linewidth} p{0.70\linewidth}}
    \toprule
    \textbf{Actor} & \textbf{Description} \\
    \midrule
    Administrator (\emph{superuser}) &
    Manages hypervisors, base templates, users, the full tree and
    global configuration. \\
    Manager &
    A professor or equivalent. Manages assigned branches of the
    tree; uses both the core and the teaching application. \\
    End user (\emph{student}) &
    Operates only the VMs assigned to them through the desk
    metaphor (DEC-009). \\
    Hypervisor &
    External system. UnivOrch talks to it through a connector;
    the hypervisor does not initiate actions. \\
    System &
    UnivOrch itself, running the operation engine and the
    automatic tasks (backup rotation, retention cleanup, failure
    detection). \\
    \bottomrule
  \end{tabular}
  \caption{Actors of UnivOrch (detailed).}
  \label{tab:actors-detail}
\end{table}

\section{Use cases}

Use cases are grouped into three families: \emph{core} use cases
(generic orchestrator operations), \emph{teaching} use cases (the
layer-2 application) and \emph{system} use cases (automatic
behaviour). The catalogue contains 32~entries; the summary below
lists every entry with its actor and one-sentence goal.

\subsection{Common}

\begin{description}
  \item[UC-AUTH-1] (Any user.) Authenticate against the system and
        obtain a session token.
\end{description}

\subsection{Core --- Administrator}

\begin{description}
  \item[UC-ADM-1] Manage users (create, edit, delete) through the
        web interface.
  \item[UC-ADM-2] Define hypervisors and their datastores with
        alias-based references.
  \item[UC-ADM-3] Define base templates (hardware parameters, base
        VM, datastore, source hypervisor).
  \item[UC-ADM-4] Create and organise folders, including optional
        IP pools that descendants inherit.
  \item[UC-ADM-5] Assign users to roles in any folder.
  \item[UC-ADM-6] Control what managers can access by deciding what
        is imported into each folder.
  \item[UC-ADM-7] Configure global properties such as Job retention.
  \item[UC-ADM-8] Force-undeploy a descriptor stuck in
        \texttt{broken} state.
  \item[UC-ADM-9] View the Job history of any branch.
  \item[UC-ADM-10] Restore the database from a backup (manual in
        v1).
  \item[UC-ADM-11] View operational reports (state of all VMs
        grouped by folder).
\end{description}

\subsection{Core --- Manager}

\begin{description}
  \item[UC-MGR-1] Create and organise folders inside the branches
        assigned to the manager, importing definitions from the
        parent.
  \item[UC-MGR-2] Create descriptors based on the templates
        allowed to the manager.
  \item[UC-MGR-3] Derive new templates from base templates; the
        manager sees the template definition but not the
        hypervisor's address or credentials (DEC-011).
  \item[UC-MGR-4] Assign students (and other managers) to roles in
        the manager's own folders.
  \item[UC-MGR-5] Trigger batch operations on a set of descriptors
        (deploy, undeploy, start, stop, pause, resume) and follow
        progress through a parent Job with child Jobs per VM.
  \item[UC-MGR-6] View the Job history restricted to the manager's
        scope.
  \item[UC-MGR-7] View a subject status report.
\end{description}

\subsection{Teaching --- Manager (layer 2)}

\begin{description}
  \item[UC-TEACH-1] Deploy a subject from a model desk and a
        student list: one desk per student, descriptors in
        \texttt{provisioned}, IPs assigned, notification emails
        sent, deployment report produced.
  \item[UC-TEACH-2] Update the student list of an existing subject;
        new students get desks, removed students are flagged
        according to policy.
  \item[UC-TEACH-3] Change a model VM; descriptors are updated and
        changes take effect on the next recreation
        (DEC-022).
\end{description}

\subsection{Core --- End user (student)}

The student sees the desk metaphor (DEC-009) and never sees the
tree.

\begin{description}
  \item[UC-USR-1] View their own desks and the state of each VM.
  \item[UC-USR-2] Deploy a VM from \texttt{provisioned}.
  \item[UC-USR-3] Undeploy a VM with a double confirmation warning,
        because the action is irreversible.
  \item[UC-USR-4] Start, stop, pause or resume a deployed VM.
  \item[UC-USR-5] Recreate a VM (undeploy followed by deploy) as
        a way to recover from a corrupted system without
        administrator help.
  \item[UC-USR-6] View the status and connection information
        (typically the IP) of a VM, read from the hypervisor on
        demand.
\end{description}

\subsection{System --- automatic behaviour}

\begin{description}
  \item[UC-SYS-1] Run the operation engine: take operations from
        the queue, execute, record the result. Synchronous in v1;
        the Job model is designed for a future asynchronous queue
        (DEC-014, DEC-015).
  \item[UC-SYS-2] Periodically back up the database following a GFS
        retention policy (DEC-024).
  \item[UC-SYS-3] Remove Jobs older than the configured retention
        period (DEC-023).
  \item[UC-SYS-4] Detect an unreachable hypervisor reactively:
        when an operation fails by communication error, the Job is
        recorded as failed and the descriptor is marked
        \texttt{unreachable}.
\end{description}

\section{Functional requirements}

The functional requirements below are grouped by area. They restate
the obligations that the use cases imply, in the form an engineer
can check off during implementation.

\subsection*{B.4.1 Tree and folders}
The system stores descriptors and folders in a hierarchical tree.
A folder has a common definition that its children import.
A parent folder \emph{publishes} definitions; the child folder
\emph{imports} what it wants to be visible below.
Folder creation is restricted by role (DEC-012).

\subsection*{B.4.2 Cascade inheritance}
The effective definition of a descriptor combines all definitions
from root to that node (DEC-010); each level may override the level
above. Hypervisors, templates, datastores and role assignments are
all inherited this way. The child folder chooses what to import
from the parent, including the wildcard~\texttt{*} (DEC-012). A
user can read the full resolved definition of every imported item,
but can only modify what is written at their own level.
Hypervisor credentials and address are never visible to anyone
below the administrator. There is no explicit \emph{ownership}
concept: it emerges from the role assignment (DEC-021).

\subsection*{B.4.3 Descriptors and states}
A descriptor holds the VM definition and the reference to the real
VM (the id used by the hypervisor) and tracks its own state. The
four descriptor states are \texttt{provisioned}, \texttt{deployed}
(with an optional \texttt{drifted} flag), \texttt{broken} and
\texttt{unreachable} (DEC-022). Runtime states are queried on
demand. Drift is detected reactively in v1.

\subsection*{B.4.4 Deployment and hypervisor connectors}
\texttt{deploy} and \texttt{undeploy} are orchestrator-level
operations; the connector exposes primitives only (DEC-016). Every
connector implements \texttt{clone}, \texttt{delete}, \texttt{start},
\texttt{stop}, \texttt{force\_stop}, \texttt{pause}, \texttt{resume},
\texttt{get\_status} and \texttt{get\_info}. \texttt{clone} takes a
mode parameter; v1 only supports \texttt{linked}. v1 ships
connectors for the mock hypervisor, VMware ESXi (designed) and
Proxmox (designed). The orchestrator is non-invasive: the
hypervisor keeps working if UnivOrch is stopped.

\subsection*{B.4.5 Jobs and the operation engine}
Every operation creates a Job with a unique identifier and the
lifecycle \texttt{pending} $\to$ \texttt{running} $\to$
\texttt{completed} or \texttt{failed} (DEC-014). Jobs cover VM
operations, tree operations and permission changes. Batch
operations spawn a parent Job and one child Job per item. Jobs are
persisted in the database (DEC-015) and provide the user-facing
history.

\subsection*{B.4.6 Users and access control}
Users are stored in a YAML file managed by the administrator from
the web interface (DEC-021). The three fixed roles are
\texttt{superuser}, \texttt{manager} and \texttt{end\_user}
(DEC-011). Role assignment lives on the folder and inherits in
cascade; it may be overridden in subfolders. The user store is
behind an abstraction so it can later be replaced (LDAP, Active
Directory).

\subsection*{B.4.7 IP address management}
Each folder may define an IP pool: a contiguous range, mask and
gateway (DEC-025). Pools cascade like every other property. At
deploy time the system picks the next free address from the
applicable pool; at undeploy time the address is released. How the
VM ultimately receives the address (DHCP, cloud-init, etc.) is
outside the scope of UnivOrch. External IPAM integration is not
included in v1 (DEC-013).

\subsection*{B.4.8 Teaching application}
The teaching application sits on top of the core and translates a
teaching workflow into core operations (DEC-004). It deploys a
subject from a model desk and a student list, applies the IP
policy, sends notification emails and produces a deployment
report. It supports updating the student list and changing the
model VM.

\subsection*{B.4.9 Reports}
v1 ships at least one operational report (VM state grouped by
folder) and one teaching report (subject status for a manager).
The full catalogue of possible reports is recorded for the future.

\subsection*{B.4.10 Interfaces}
The system exposes a REST API used by external clients (DEC-018).
A CLI works as both single commands and an interactive shell. A
web interface covers all roles and is designed especially for the
student. The CLI is a pure REST client of the daemon. The web
interface is mounted on the daemon's HTTP server.

\subsection*{B.4.11 Persistence, logging and backup}
Tree, descriptors, Jobs and users are persisted and survive
restarts. Persistence goes through the Repository pattern (DEC-007).
System logs use Python's standard \texttt{logging} (DEC-023).
Operation logs are the Job history with configurable retention
(default 90~days). The database is backed up periodically under a
GFS retention policy (DEC-024); v1 restores are manual.

\section{Non-functional requirements}

\Cref{tab:nfr-detail} restates the non-functional requirements as
soft targets that keep the implementation honest, not as
contractual service levels.

\begin{table}[ht]
  \centering
  \small
  \begin{tabular}{p{0.22\linewidth} p{0.70\linewidth}}
    \toprule
    \textbf{Quality} & \textbf{Target} \\
    \midrule
    Usability     & A non-technical user operates their VMs without
                    any knowledge of hypervisors. \\
    Security      & Every operation requires authentication. Users
                    only see and act on resources allowed by their
                    role and folder assignments. Plain-text password
                    storage is a known PoC limitation. \\
    Performance   & A realistic teaching load: hundreds of VMs
                    across several subjects. Cascade inheritance
                    keeps the configuration manageable at that
                    scale. \\
    Reliability   & The system stays consistent across restarts. A
                    failed operation always leaves a record in the
                    Job history. An inconsistent descriptor goes to
                    \texttt{broken}. The database can be restored
                    from a backup. \\
    Maintainability & New hypervisors are added by implementing
                    the common connector interface, without touching
                    the core. New applications are built on the
                    core without modifying it (DEC-004). \\
    Portability   & The system runs as a Linux service. It does
                    not require the hypervisors on the same host
                    and does not interfere with their normal
                    operation. \\
    \bottomrule
  \end{tabular}
  \caption{Non-functional requirements of UnivOrch v1.}
  \label{tab:nfr-detail}
\end{table}

\section{Constraints and assumptions}

\begin{itemize}
  \item The implementation language is Python; the version is
        chosen in Phase~4 (DEC-033).
  \item The system runs as a service on Linux.
  \item Hypervisors exist and are not part of UnivOrch; base VMs
        are created directly on the hypervisor by the
        administrator.
  \item Operations are synchronous in v1.
  \item Passwords are stored in plain text in v1 (PoC only).
  \item IP assignment uses simple per-folder pools in v1; external
        IPAM integration is out of scope (DEC-013, DEC-025).
\end{itemize}

\section{Out of scope for v1}

\begin{itemize}
  \item Automatic reconciliation loop (changes are applied on
        demand) -- DEC-006.
  \item High availability and clustering.
  \item External IPAM integration -- DEC-013.
  \item Snapshot management through the orchestrator.
  \item Reference-only descriptors and discovery of pre-existing
        VMs -- DEC-005b, DEC-020.
  \item LDAP / Active Directory integration -- DEC-021.
  \item Text-mode user interface (TUI) -- DEC-018.
  \item Web interface for backup restore -- DEC-024.
  \item Per-folder fine-grained permission customisation --
        DEC-011.
\end{itemize}
